\documentclass{article}
\usepackage[margin=1.0in]{geometry}
\usepackage{booktabs}
\usepackage{makecell}
\usepackage{float}
\usepackage{caption}
\usepackage{amsmath,amssymb,amsfonts}
\usepackage{dsfont}
\usepackage[normalem]{ulem}
\usepackage{setspace}
\setstretch{1.05}

\newcommand{\cleanedresultsdir}{../../results/cleaned/tex}
\newcommand{\TableInput}[1]{\input{#1}}
\DeclareCaptionStyle{noteflush}{%
  format=plain,
  justification=raggedright,
  singlelinecheck=false,
  font=footnotesize,
  width=\linewidth,
  skip=0pt
}
\newcommand{\FloatNote}[1]{%
  \begingroup
    \captionsetup{style=noteflush}%
    \caption*{\textit{Notes:}~#1}%
  \endgroup
}

\begin{document}

\section*{User Productivity: Startup-swapped horse races and equity specs}

\noindent This note reports two user-productivity exercises built off the current core design. The first keeps the baseline startup triple-difference in place and re-runs the horse race with startup-swapped controls on the remote margin. The second estimates a separate equity quadruple that appends the equity signal directly to the core empirical setup.

\section*{1. Startup-swapped horse race}

\subsection*{Empirical setup}
\noindent For each added control $Z_f$ in the current horse-race stack, the modified specification keeps the baseline startup term and adds a startup-swapped remote channel:
\[
  \begin{aligned}
    y_{ifh} =\; & \beta_3 \left( \text{Remote}_f \times \mathds{1}(\text{Post})_h \right)
    + \beta_5 \left( \text{Remote}_f \times \mathds{1}(\text{Post})_h \times \text{Startup}_f \right) \\
    & + \sum_{Z \in \mathcal{Z}} \left[
        \gamma_Z \left( \text{Remote}_f \times \mathds{1}(\text{Post})_h \times Z_f \right)
        + \theta_Z \left( \mathds{1}(\text{Post})_h \times Z_f \right)
      \right]
    + \text{FE} + \varepsilon_{ifh} .
  \end{aligned}
\]
\noindent Here $\mathcal{Z}$ contains the two competing channels in the note: a firm-level bin for average post-COVID employment growth and the firm-level equity-offer signal. In IV, each remote interaction $\text{Remote}\times\mathds{1}(\text{Post})\times Z_f$ is instrumented by the matching $\text{Teleworkable}\times\mathds{1}(\text{Post})\times Z_f$ interaction. Within each fixed-effects family, every column is estimated on the same merged sample.

\subsection*{Baseline FE}
\begin{table}[H]
  \centering
  \caption{User productivity startup-swapped horse race, baseline FE}
  \TableInput{\cleanedresultsdir/llm_equity_startup_swapped_horse_race_baseline.tex}
  \FloatNote{Columns mirror the current baseline-FE horse-race structure: baseline, firm growth only, equity only, and both. The first two rows report the baseline remote coefficients on $\text{Remote}\times\mathds{1}(\text{Post})$ and $\text{Remote}\times\mathds{1}(\text{Post})\times\text{Startup}$. The growth channel uses a firm-level median split of average post-COVID employment growth. The added rows use startup-swapped controls: $\mathds{1}(\text{Post})\times Z$ and $\text{Remote}\times\mathds{1}(\text{Post})\times Z$ for firm growth and equity, with the remote interactions instrumented by the matching teleworkability interactions. Standard errors are clustered by worker.}
\end{table}

\subsection*{Pair FE}
\begin{table}[H]
  \centering
  \caption{User productivity startup-swapped horse race, pair FE}
  \TableInput{\cleanedresultsdir/llm_equity_startup_swapped_horse_race_pairfe.tex}
  \FloatNote{Columns mirror the current pair-FE horse-race structure: baseline, firm growth only, equity only, and both. The growth channel uses the same firm-level median split of average post-COVID employment growth as in Table 1. The added rows use the same startup-swapped controls as in Table 1, estimated with firm$\times$individual and half-year fixed effects. Standard errors are clustered by worker.}
\end{table}

\section*{2. Equity quadruple}

\subsection*{Empirical setup}
\noindent The separate equity-quadruple exercise appends the equity signal directly to the core empirical setup:
\[
  \begin{aligned}
    y_{ifh} =\; & \beta_1 \left( \text{Remote}_f \times \mathds{1}(\text{Post})_h \right)
    + \beta_2 \left( \mathds{1}(\text{Post})_h \times \text{Startup}_f \right)
    + \beta_3 \left( \text{Remote}_f \times \mathds{1}(\text{Post})_h \times \text{Startup}_f \right) \\
    & + \delta_1 \left( \text{Remote}_f \times \mathds{1}(\text{Post})_h \times \text{Offers equity}_f \right) \\
    & + \delta_2 \left( \mathds{1}(\text{Post})_h \times \text{Startup}_f \times \text{Offers equity}_f \right) \\
    & + \delta_3 \left( \text{Remote}_f \times \mathds{1}(\text{Post})_h \times \text{Startup}_f \times \text{Offers equity}_f \right)
    + \text{FE} + \varepsilon_{ifh} .
  \end{aligned}
\]
\noindent In IV, the remote-post term, the remote-post-startup term, the remote-post-equity term, and the remote-post-startup-equity term are instrumented with the matching teleworkability interactions. The post-startup term and the post-startup-equity term enter as included exogenous controls.

\subsection*{Baseline FE}
\begin{table}[H]
  \centering
  \caption{User productivity equity quadruple, baseline FE}
  \TableInput{\cleanedresultsdir/llm_equity_quadruple_baseline_fe.tex}
  \FloatNote{Column (1) is the baseline user-productivity specification. Column (2) appends the equity signal to the core empirical setup through $\text{Remote}\times\mathds{1}(\text{Post})\times\text{Offers equity}$, $\mathds{1}(\text{Post})\times\text{Startup}\times\text{Offers equity}$, and $\text{Remote}\times\mathds{1}(\text{Post})\times\text{Startup}\times\text{Offers equity}$. The remote interactions are instrumented by the matching teleworkability interactions. Standard errors are clustered by worker.}
\end{table}

\subsection*{Pair FE}
\begin{table}[H]
  \centering
  \caption{User productivity equity quadruple, pair FE}
  \TableInput{\cleanedresultsdir/llm_equity_quadruple_pair_fe.tex}
  \FloatNote{Column (1) is the pair-FE baseline specification. Column (2) appends the equity signal to the core empirical setup under firm$\times$individual and half-year fixed effects. The remote interactions are instrumented by the matching teleworkability interactions. Standard errors are clustered by worker.}
\end{table}

\end{document}
